\documentclass{beamer}
%\documentclass[handout]{beamer}
\mode<presentation>
{
%  \usetheme{Darmstadt}%default Warsaw
  \setbeamercovered{transparent}
}

\usepackage{beamerthemesplit}
\usepackage{epsfig,graphics}
\usepackage{amsmath,amsfonts,amssymb,color,array,subfigure,rotating,texdraw} 
\usepackage{multimedia}
\usepackage{psfrag}
\usepackage{subfigure}
%\usepackage{unicode-math}
%\usepackage[dvips]{color}

\newcommand{\bDiamond}{\mathbin{\Diamond}}
\newcommand{\bLozenge}{\mathbin{\blacklozenge}}

\makeatletter
\newcommand\bigDiamond{\mathop{\mathpalette\bigDi@mond\relax}}
\newcommand\bigDi@mond[2]{%
  \vcenter{\hbox{\m@th
    \scalebox{\ifx#1\displaystyle 2\else1.2\fi}{$#1\Diamond$}%
  }}%
}
\newcommand\bigLozenge{\mathop{\mathpalette\bigL@zenge\relax}}
\newcommand\bigL@zenge[2]{%
  \vcenter{\hbox{\m@th
    \scalebox{\ifx#1\displaystyle 2\else1.2\fi}{$#1\blacklozenge$}%
  }}%
}
\makeatother



\definecolor{greenn}{rgb}{0.1,0.4,0.1}
\def\Dpartial#1#2{ {\partial #1 \over \partial #2} }
\def\Dpartialmix#1#2#3{ {\partial^2 #1 \over \partial #2\,\partial #3} }
\def\Dpartialn#1#2#3{ {\partial^{#3} #1 \over \partial #2^{#3}} }
\def\Txtint{ {\textstyle\int} }
\def\Txtiint{ {\textstyle\iint} }
\def\Bmp#1{ \begin{minipage}{#1} }
\def\Emp{ \end{minipage} }
\def\Bmpc#1{ \begin{minipage}[c]{#1} }
\def\Bmpt#1{ \begin{minipage}[t]{#1} }
\def\Bmpb#1{ \begin{minipage}[b]{#1} }
\newcommand{\stitle}[1]{\centering{{\sc \Large #1 }}\vskip .5cm}
\def\H{{\mathcal{H}}}
\def\J{{\mathcal{J}}}

\def\A{{\mathcal{A}}}
\def\B{{\mathcal{B}}}
\def\K{{\mathcal{K}}}
\def\L{{\mathcal{L}}}
\def\M{{\mathcal{M}}}
\def\O{{\mathcal{O}}}
\def\F{{\mathcal{F}}}
\def\N{{\mathcal{N}}}
\def\T{{\mathcal{T}}}
\def\R{{\mathcal{R}}}
\def\G{{\mathcal{G}}}
\def\E{{\mathcal{E}}}
\def\K{{\mathcal{K}}}
\def\C{{\mathcal{C}}}
\def\U{{\mathcal{U}}}
\def\S{{\mathcal{S}}}
\def\DD{{\mathcal{D}}}
\def\QQ{{\mathcal{Q}}}
\def\NS{{\mathcal{NS}}}
\def\l{{\ell}}

\def\BB{{\mathbb{B}}}
\def\CC{{\mathbb{C}}}
\def\KK{{\mathbb{K}}}
\def\RR{{\mathbb{R}}}
\def\TT{{\mathbb{T}}}

\def\tu{{\tilde{u}}}
\def\ta{{\tilde{a}}}
\def\tb{{\tilde{b}}}

\def\magenta#1{{\color{magenta}{ #1 }}}
\def\red#1{{\color{red}{ #1 }}}
\def\blue#1{{\color{blue}{ #1 }}}
\def\green#1{{\color{green}{ #1 }}}
\def\white#1{{\color{white}{ #1 }}}
\def\black#1{{\color{black}{ #1 }}}

\def\bA{{\bf A}}
\def\bB{{\bf B}}
\def\bC{{\bf C}}
\def\bF{{\bf F}}
\def\bF{{\bf F}}
\def\bK{{\bf K}}
\def\bM{{\bf M}}
\def\bZ{{\bf Z}}


\def\P{{\mathcal{P}}}
\def\b{{\bf b}}
\def\q{{\bf q}}
\def\r{{\bf r}}
\def\v{{\bf v}}
\def\z{{\bf z}}
\def\y{{\bf y}}
\def\k{{\bf k}}
\def\x{{\bf x}}
\def\m{{\bf m}}
\def\l{{\bf l}}
\def\n{{\bf n}}
\def\t{{\bf t}}
\def\e{{\bf e}}
\def\D{{\bf D}}
\def\X{{\bf X}}
\def\Q{{\bf Q}}
\def\g{{\bf g}}
\def\f{{\bf f}}
\def\u{{\bf u}}
\def\y{{\bf y}}
\def\0{{\mathbf{0}}}
\def\V{{\mathbf{V}}}
\def\W{{\mathbf{W}}}
\def\Z{{\mathbf{Z}}}

\def\tT{{\tilde{\mathcal{T}}}}
\def\tZ{{\tilde{\mathcal{Z}}}}
\def\tM{{\tilde{M}}}

\newcommand{\bds}{\boldsymbol}
\newcommand{\mbb}[1]{\mathbb{#1}}
\newcommand{\mc}[1]{\mathcal{#1}}


\def\btau{{\bf \tau}}
\def\hu{{\hat{u}}}
\def\bnabla{\boldsymbol{\nabla}}
\def\bphi{\boldsymbol{\phi}}
\def\bDelta{\boldsymbol{\Delta}}
\def\bPhi{\boldsymbol{\Phi}}
\def\bxi{\boldsymbol{\xi}}
\def\etab{\boldsymbol{\eta}}
\def\control{{\phi}}
\def\controlpert{{\phi'}}
\def\bsigma{\boldsymbol{\sigma}}
\def\tomega{{\widetilde{\omega}}}
\def\oomega{{\overline{\omega}}}
\def\tpsi{{\widetilde{\psi}}}

\newcommand{\sigdis}{\sigma_{\text{disc}}}
\newcommand{\sigess}{\sigma_{\text{ess}}}
\def\sigp{\sigma_{\text{p}}}
\def\mumax{\mu_{\text{max}}}

\newcommand{\argmin}{\operatorname{argmin}}
\newcommand{\argmax}{\operatorname{argmax}}
\newcommand{\sgn}{\operatorname{sgn}}
\newcommand{\sign}{\operatorname{sign}}
\newcommand{\rank}{\operatorname{rank}}
\newcommand{\Ext}{\operatorname{Ext}}
\newcommand{\diag}{\operatorname{diag}}
\newcommand{\Range}{\operatorname{Range}}

%%%%%%%%%%%%%%%%%%%%%%%%%%%%%%%%%%%%%%%%%%%%%%%%%%%%%%%%%%%%%%%%%%%%%%%%%%
%%%%%%%%%%%%%%%%%%%%%%%%%%%%%%%%%%%%%%%%%%%%%%%%%%%%%%%%%%%%%%%%%%%%%%%%%%
%%%%%%%%%%%%%%%%%%%%%%%%%%%%%%%%%%%%%%%%%%%%%%%%%%%%%%%%%%%%%%%%%%%%%%%%%%

\title[On Instabilities of the Taylor-Green Vortex in 2D Euler Flows]
{On Instabilities of the Taylor-Green Vortex \\ in 2D Euler Flows}
%\subtitle{Extreme Vortex States and \\ the Hydrodynamic Blow-Up Problem}
\author[X.~Zhao, \underline{B.~Protas} and R.~Shvydkoy]{Xinyu Zhao$^{1,2}$, \underline{Bartosz Protas}$^{2}$ and Roman Shvydkoy$^{3}$}
\institute{
$^{1}$Department of Mathematics \& Statistics \\
McMaster University, Hamilton, Ontario, Canada \\ \vspace*{0.3cm}
$^2$Department of Mathematical Sciences, New Jersey Institute of Technology, \\
Newark, NJ 07103, USA \\ \vspace*{0.3cm}
$^3$Department of Mathematics, Statistics and Computer Science, \\
University of Illinois at Chicago, Chicago, IL 60607, USA \\ \vspace*{0.25cm}
%\footnotesize
Funded by NSERC (Canada), Computing Time Provided by {\sc Compute Canada} \vspace*{-0.25cm} }

\date{\scriptsize %\vspace*{-1.0cm}
 Mathematical Congress of the Americas, \quad Miami, FL \quad July 12--25, 2025 \\
 Session {\it "New Developments in Mathematical Fluid Dynamics"}}




\begin{document}

\frame{\titlepage}


%%%%%%%%%%%%%%%%%%%%%%%%%%%%%%%%%%%%%%%%%%%%%%%%%%%%%%%%%%%%%%%%%%%%%%%%%%
%%%%%%%%%%%%%%%%%%%%%%%%%%%%%%%%%%%%%%%%%%%%%%%%%%%%%%%%%%%%%%%%%%%%%%%%%%
%%%%%%%%%%%%%%%%%%%%%%%%%%%%%%%%%%%%%%%%%%%%%%%%%%%%%%%%%%%%%%%%%%%%%%%%%%
\frame{
%\frametitle{References}
\centering
\medskip
\includegraphics[width=1.0\textwidth]{Figs9/zps24a2.png}
\bigskip\smallskip
}


%%%%%%%%%%%%%%%%%%%%%%%%%%%%%%%%%%%%%%%%%%%%%%%%%%%%%%%%%%%%%%%%%%%%%%%%%%
%%%%%%%%%%%%%%%%%%%%%%%%%%%%%%%%%%%%%%%%%%%%%%%%%%%%%%%%%%%%%%%%%%%%%%%%%%
%%%%%%%%%%%%%%%%%%%%%%%%%%%%%%%%%%%%%%%%%%%%%%%%%%%%%%%%%%%%%%%%%%%%%%%%%%

\frame
{
\hspace*{-1.0cm}
\Bmp{12.5cm}
\raggedright
\frametitle{Vortex Rings in Nature and Technology}

\centering
\includegraphics[width=0.9\textwidth]{Figs12/vringX}
\Emp
}


%%%%%%%%%%%%%%%%%%%%%%%%%%%%%%%%%%%%%%%%%%%%%%%%%%%%%%%%%%%%%%%%%%%%%%%%%%
%%%%%%%%%%%%%%%%%%%%%%%%%%%%%%%%%%%%%%%%%%%%%%%%%%%%%%%%%%%%%%%%%%%%%%%%%%
%%%%%%%%%%%%%%%%%%%%%%%%%%%%%%%%%%%%%%%%%%%%%%%%%%%%%%%%%%%%%%%%%%%%%%%%%%

%\section[Agenda]{}
\frame{
\frametitle{Agenda}
\tableofcontents
}



%%%%%%%%%%%%%%%%%%%%%%%%%%%%%%%%%%%%%%%%%%%%%%%%%%%%%%%%%%%%%%%%%%%%%%%%%%
%%%%%%%%%%%%%%%%%%%%%%%%%%%%%%%%%%%%%%%%%%%%%%%%%%%%%%%%%%%%%%%%%%%%%%%%%%
%%%%%%%%%%%%%%%%%%%%%%%%%%%%%%%%%%%%%%%%%%%%%%%%%%%%%%%%%%%%%%%%%%%%%%%%%%
\section{Introduction}
\subsection{Taylor-Green Vortex}
\frame
{
\hspace*{-0.75cm}
\Bmp{12cm}
\raggedright
\vspace*{0.1cm}
\begin{itemize}
\item<1-> 2D Euler system on a periodic domain
\begin{equation*}
\begin{aligned}
\partial_t \omega + (\bds u \cdot \nabla)\omega = 0&, &\qquad &(\bds x, t) \in 
 \mbb T^2 \times (0, T], \\
 \omega|_{t = 0}  = \omega_0 &, &\qquad &\bds x \in \mbb T^2,
\end{aligned}
\end{equation*}
where 
\begin{itemize}
\item<2-> $\blue{\omega  = \nabla\times \bds u}$ is the scalar vorticity field satisfying 
$\int_{\mbb T^2} \omega \; d\bds x = 0$
\smallskip
\item<2-> Biot-Savart law: $\bds u = \nabla^\perp \Delta^{-1} \omega, \qquad \nabla^\perp = (-\partial_y, \partial_x)$
\end{itemize}

\medskip
\item<3-> An equilibrium solution --- the 2D Taylor-Green vortex: 

\mbox{
\Bmpb{0.5\textwidth}
%\bigskip
\vspace*{-1.5cm}
\blue{
\begin{align*} 
\bds u_s & = (-\cos(x_1)\sin(x_2), \sin(x_1)\cos(x_2)) \\ \\
\omega_s & = 2\cos(x_1)\cos(x_2)
\end{align*}}
\Emp
\Bmpb{0.5\textwidth}
\centering
\visible<3->{\includegraphics[width=0.8\textwidth] {./Figs9/TaylorGreen.png}}
\Emp} 

\end{itemize}

     
\Emp
}
%%%%%%%%%%%%%%%%%%%%%%%%%%%%%%%%%%%%%%%%%%%%%%%%%%%%%%%%%%%%%%%%%%%%%%%%%%
%%%%%%%%%%%%%%%%%%%%%%%%%%%%%%%%%%%%%%%%%%%%%%%%%%%%%%%%%%%%%%%%%%%%%%%%%%
%%%%%%%%%%%%%%%%%%%%%%%%%%%%%%%%%%%%%%%%%%%%%%%%%%%%%%%%%%%%%%%%%%%%%%%%%%
\subsection{Linearized Euler Operator}
\frame
{
\hspace*{-0.75cm}
\Bmp{12cm}
\raggedright

\begin{itemize}
\item<1-> Linearize the 2D Euler system around a steady solution $(\omega_s, \bds u_s)$ $\Longrightarrow$
\begin{align*}
\partial_t w &= {\blue \L} w, \qquad w|_{t = 0} = w_0, \qquad \text{where} \\
{\color {blue} \L w} :&= -(\bds u_s\cdot \nabla) w -(\bds u \cdot \nabla) \omega_s \\
& = {\color{blue} \left[-\bds u_s \cdot \nabla - \nabla\omega_s\cdot (\nabla^\perp \Delta^{-2})\right] w}, \qquad w \in H^m_0(\mbb T^2)
\end{align*}

%\smallskip
\item<2-> Linear stability of the equilibria determined by the spectrum of the operator $\L$
\blue{
\begin{align*}
\sigma(\L) := & \left\{ \lambda \in \CC \; : \; \text{$(\L - \lambda)$ is not bounded from below in $\mathcal{X}$} \right\} \\ 
= & \ \sigma_{\text{disc}}(\L) \cup \sigma_{\text{ess}} (\L)
\end{align*}}
\small \vspace*{-0.4cm}
\begin{itemize}
\item<3-> A point $z  \in \sigdis(\L)$ if
\begin{itemize}
\item<3->[(1)] $z$ is an isolated point in $\sigma(\L)$;
\item<3->[(2)]$z$ has finite multiplicity, i.e., $\bigcup_{r=1}^\infty
    \mathrm{Ker}(z - \L)^r$ is finite dimensional;
\item<3->[(3)] The range of $z-\L$ is closed. 
\end{itemize}
\smallskip
\item<4-> Otherwise, $z  \in \sigess(\L)$
\end{itemize}

\end{itemize}

     
\Emp
}

%%%%%%%%%%%%%%%%%%%%%%%%%%%%%%%%%%%%%%%%%%%%%%%%%%%%%%%%%%%%%%%%%%%%%%%%%%
%%%%%%%%%%%%%%%%%%%%%%%%%%%%%%%%%%%%%%%%%%%%%%%%%%%%%%%%%%%%%%%%%%%%%%%%%%
%%%%%%%%%%%%%%%%%%%%%%%%%%%%%%%%%%%%%%%%%%%%%%%%%%%%%%%%%%%%%%%%%%%%%%%%%%
\subsection{Spectrum}
\frame
{
\hspace*{-0.75cm}
\Bmp{12cm}
\raggedright

\begin{itemize}
\item<1-> 
The set of eigenvalues of $\L$ forms its point spectrum
\begin{equation*}
\blue{{\sigp(\L) := \left\{ \lambda \in \CC \; : \; \exists \phi (\x) \neq 0,  \quad \L \phi (\x) = \lambda \phi (\x), \quad \x \in \TT^2  \right\}}}
\end{equation*}

\bigskip
\item<2->
In finite dimensions linear operators can only have point spectrum

\bigskip
\item<3-> In infinite dimensions the situation is complicated by the
  presence of the essential spectrum $\sigess(\L)$, which may also contain
  eigenvalues of $\L$ that are not in $\sigdis(\L)$

\end{itemize}

     
\Emp
}


%%%%%%%%%%%%%%%%%%%%%%%%%%%%%%%%%%%%%%%%%%%%%%%%%%%%%%%%%%%%%%%%%%%%%%%%%%
%%%%%%%%%%%%%%%%%%%%%%%%%%%%%%%%%%%%%%%%%%%%%%%%%%%%%%%%%%%%%%%%%%%%%%%%%%
%%%%%%%%%%%%%%%%%%%%%%%%%%%%%%%%%%%%%%%%%%%%%%%%%%%%%%%%%%%%%%%%%%%%%%%%%%

\frame
{
\hspace*{-0.75cm}
\Bmp{12cm}
\raggedright

\begin{itemize}
\item<1-> For the 2D linear Euler operator $\L$ the essential spectrum
  $\sigess(\L)$ is well understood (Friedlander et al.~1997, Shvydkoy
  \& Latushkin 2003)
\begin{equation*}
  \blue{\sigess(\L; H^{m}_0) = \{ z\in \CC: |\Re(z)| \leq |m| \, \mumax \}},
\end{equation*}
where 
\begin{equation*}
\mumax :=  \lim_{t\to\infty} \frac{1}{t}\log \sup_{\bds x\in\mbb T^2} ||\nabla \varphi_t(\bds x)||
\end{equation*}
is the maximal Lyapunov exponent corresponding
to the Lagrangian flow $\varphi_t: \bxi \to \bds x(t; \bxi)$ generated
by the steady state via $\partial_t \bds x(t) = \bds u_s (\bds x(t))$

\medskip
\item<2-> If $z \in \sigess(\L)$, then there exists a sequence of
  weak-null unit vectors
  ${\{f_n\} \in H^m_0(\TT^2)}$,\/\blue{approximate eigenvectors,}\/such
  that
  \blue{$$\|(\L - z) f_n \|_{H^m} \rightarrow 0, \qquad n \rightarrow
    \infty$$} and $\{f_n\}$ does not contain any convergent
  subsequence
\end{itemize}

     
\Emp
}

%%%%%%%%%%%%%%%%%%%%%%%%%%%%%%%%%%%%%%%%%%%%%%%%%%%%%%%%%%%%%%%%%%%%%%%%%%
%%%%%%%%%%%%%%%%%%%%%%%%%%%%%%%%%%%%%%%%%%%%%%%%%%%%%%%%%%%%%%%%%%%%%%%%%%
%%%%%%%%%%%%%%%%%%%%%%%%%%%%%%%%%%%%%%%%%%%%%%%%%%%%%%%%%%%%%%%%%%%%%%%%%%

\frame
{
\hspace*{-0.75cm}
\Bmp{12cm}
\raggedright

\begin{itemize}
\item<1->
The the maximal Lyapunov exponent is determined by the hyperbolic stagnation point 
\begin{equation*}
\mumax = \max\left[ \lambda\left(\bnabla\u_s({\x_s}) \right)\right]
\end{equation*}

\medskip
\item<2->
  In the Taylor-Green vortex: \qquad \red{$\mumax = 1$} \\
  \medskip
  
  \centering
\visible<2->{\includegraphics[width=0.6\textwidth] {./Figs9/TaylorGreen.png}}

\end{itemize}

     
\Emp
}



%%%%%%%%%%%%%%%%%%%%%%%%%%%%%%%%%%%%%%%%%%%%%%%%%%%%%%%%%%%%%%%%%%%%%%%%%%
%%%%%%%%%%%%%%%%%%%%%%%%%%%%%%%%%%%%%%%%%%%%%%%%%%%%%%%%%%%%%%%%%%%%%%%%%%
%%%%%%%%%%%%%%%%%%%%%%%%%%%%%%%%%%%%%%%%%%%%%%%%%%%%%%%%%%%%%%%%%%%%%%%%%%

\frame
{
\hspace*{-0.75cm}
\Bmp{12cm}
\raggedright

\begin{itemize}
\item<1-> No general results available about the discrete spectrum $\sigdis(\L)$

\smallskip
\item<2-> Unstable eigenvalues found for shear flows (Drazin \& Reid
  1981; Friedlander et al.~1997), Kolmogorov flows (Friedlander et
  al.~1997) and the cellular ``cat’s-eye'' flow (Friedlander et
  al.~2000).

\medskip
\item<3->
  The problem remains open for the Taylor-Green vortex
  \smallskip
\begin{itemize}
\item<4-> Some work in the viscous setting (Sipp \&
  Jacquin 1998; Leblanc \& Godeferd 1999; Hattori \& Hirota
  2023). However, in that case the vortex decays as
  $\mathcal{O}(e^{-\nu t})$ and\/\blue{$\sigess(\L) = \emptyset$}
\end{itemize}

\smallskip
\item<5->
{\red{\sc Question:}} \ \blue{does $\L$ have unstable eigenvalues and, if
  so, do they belong to $\sigdis (\L)$ or $\sigess (\L)$?} \\ \vspace*{-0.3cm}
  \Bmp{\textwidth}
  \hspace*{5.0cm} 
\visible<5->{\includegraphics[width=0.35\textwidth]{Figs12/piess_01.pdf}}
\Emp
\end{itemize}

     
\Emp
}



%%%%%%%%%%%%%%%%%%%%%%%%%%%%%%%%%%%%%%%%%%%%%%%%%%%%%%%%%%%%%%%%%%%%%%%%%%
%%%%%%%%%%%%%%%%%%%%%%%%%%%%%%%%%%%%%%%%%%%%%%%%%%%%%%%%%%%%%%%%%%%%%%%%%%
%%%%%%%%%%%%%%%%%%%%%%%%%%%%%%%%%%%%%%%%%%%%%%%%%%%%%%%%%%%%%%%%%%%%%%%%%%

\frame
{
\hspace*{-0.75cm}
\Bmp{12cm}
\raggedright
\footnotesize
\begin{theorem}[Lin 2004]
  Suppose there exists an exponentially growing solution
  $e^{\lambda t} w_0$ of the linearized system $\partial_t w = \L w$
  with $\Re(\lambda) > 0$ and let\/\blue{$w_0 \in L^2(\TT^2)$.} Then we have
  the following
\begin{itemize}
\item[(i)] [regularity of growing modes] \red{$w_0 \in W^{1,p}(\TT^2)
  \bigcap L^q(\TT^2)$} for all $1 \le p < p^*$ and $1 \le q < \infty$,
  where
  \begin{equation*}
    p^* = \begin{cases}
      \red{\frac{\mumax}{\mumax - \Re(\lambda)}} = \frac{1}{1 -\Re(\lambda)}, \qquad & \mumax > \Re(\lambda) \\
      \infty, & \mumax \le \Re(\lambda)
      \end{cases}
    \end{equation*}

  \item[(ii)] [nonlinear instability] if there exists $\epsilon > 0$,
    such that for any $\delta > 0$ there is a solution $w^\delta(t)$
    of the 2D Euler system corresponding to the initial condition
    $\omega_0^{\delta}$, such that
\begin{equation*}
 \| \omega_0^{\delta} - \omega_s \|_{L^q} +  \| \omega_0^{\delta} - \omega_s \|_{W^{1,p}} \le \delta,
\end{equation*}
and
\begin{equation*}
{\sup_{0 < t < T_\delta} \|\omega^\delta(t) - \omega_s\|_{H^m} \geq \epsilon,}
\end{equation*}
%
where $T_\delta = O(|\ln(\delta)|)$, $q \in [1, \infty)$, $p =
[1,p^*)$ and $m \ge -1$.
\end{itemize}
\end{theorem}

     
\Emp
}

%%%%%%%%%%%%%%%%%%%%%%%%%%%%%%%%%%%%%%%%%%%%%%%%%%%%%%%%%%%%%%%%%%%%%%%%%%
%%%%%%%%%%%%%%%%%%%%%%%%%%%%%%%%%%%%%%%%%%%%%%%%%%%%%%%%%%%%%%%%%%%%%%%%%%
%%%%%%%%%%%%%%%%%%%%%%%%%%%%%%%%%%%%%%%%%%%%%%%%%%%%%%%%%%%%%%%%%%%%%%%%%%
\section{Modal Growth of Perturbations}
\subsection{Unstable Eigenvalues and Eigenvectors}
\frame
{
\hspace*{-0.75cm}
\Bmp{12cm}
\raggedright

\begin{itemize}
\item<1-> Numerical solution of the eigenvalue problem \blue{$\L \phi
    (\x) = \lambda \phi (\x)$}

 \medskip
\item<2->  Discretize $\L$ using the basis
  \footnotesize
\begin{equation*}
\begin{gathered}
\varphi_{j_1, j_2} = \frac{1}{\sqrt{2} \pi}\cos(j_1x + j_2y), \qquad 
\psi_{j_1, j_2} = -\frac{1}{\sqrt{2} \pi}\sin(j_1x + j_2y), \\
\qquad (j_1, j_2) \in \{(0, j_2), 1\leq j_2 \leq N\} \bigcup 
\{(j_1, j_2), 1\leq j_1 \leq N, -N\leq j_2\leq N\}
\end{gathered}
\end{equation*}
\normalsize
to obtain a sparse algebraic eigenvalue problem  \quad \blue{$\mathbf{L} \,
  \phi = \lambda \phi$}

\medskip
\item<3-> The spectrum of $\mathbf{L}$
\mbox{
  \visible<3->{\includegraphics[width=0.4\textwidth]{./Figs9/spectrum_TG}} \qquad
  \visible<4->{\includegraphics[width=0.4\textwidth]{./Figs9/eigs_Krylov}}}
\end{itemize}

     
\Emp
}


%%%%%%%%%%%%%%%%%%%%%%%%%%%%%%%%%%%%%%%%%%%%%%%%%%%%%%%%%%%%%%%%%%%%%%%%%%
%%%%%%%%%%%%%%%%%%%%%%%%%%%%%%%%%%%%%%%%%%%%%%%%%%%%%%%%%%%%%%%%%%%%%%%%%%
%%%%%%%%%%%%%%%%%%%%%%%%%%%%%%%%%%%%%%%%%%%%%%%%%%%%%%%%%%%%%%%%%%%%%%%%%%

\frame
{
\hspace*{-0.75cm}
\Bmp{12cm}
\raggedright

\begin{itemize}
\item<1-> The unstable eigenvalue: \quad \red{$\lambda_+^{3000} =
    0.1424 + 0.5875i$}

  \medskip
\item<2-> The corresponding eigenvectors $\Re\left[\phi_+^N\right]$ are
  $L^2$ integrable
\visible<2->{\mbox{
  \includegraphics[width=0.425\textwidth]{./Figs9/eigfun3000} \quad
  \includegraphics[width=0.425\textwidth]{./Figs9/vort_section}}
}

% \medskip
\vspace*{-0.25cm}
\item<3-> From the theorem of Lin (2004):
  \small \vspace*{-0.25cm}
  \begin{equation*}
\phi_+ \in W^{1,p^*}(\TT^2), \qquad \text{where} \qquad \blue{p^* = \frac{1}{1 -
    0.14} = 1.16}
\end{equation*}

\normalsize
\item<3-> From the Sobolev embedding  $W^{1,p^*} \hookrightarrow H^s$ with
  $1/p^* - 1/2 = 1/2 - s / 2$, we have
    \small \vspace*{-0.35cm}
  \begin{equation*}
\qquad \qquad \qquad \qquad \blue{\phi_+ \in H^{0.28}(\TT^2)}
\end{equation*}
  
\end{itemize}

     
\Emp
}

%%%%%%%%%%%%%%%%%%%%%%%%%%%%%%%%%%%%%%%%%%%%%%%%%%%%%%%%%%%%%%%%%%%%%%%%%%
%%%%%%%%%%%%%%%%%%%%%%%%%%%%%%%%%%%%%%%%%%%%%%%%%%%%%%%%%%%%%%%%%%%%%%%%%%
%%%%%%%%%%%%%%%%%%%%%%%%%%%%%%%%%%%%%%%%%%%%%%%%%%%%%%%%%%%%%%%%%%%%%%%%%%

\frame
{
\hspace*{-0.75cm}
\Bmp{12cm}
\raggedright

The Fourier spectrum of the eigenfunction $\phi$
\\ \vspace*{-0.4cm}
\begin{equation*}
e(k) := \frac{1}{2}\sum_{k \leq |\bds j| < k+1} |\widehat\phi_{\bds j}|^2, \qquad k \in \mbb N
\end{equation*}
\vspace*{-0.3cm}

\centering
\includegraphics[width=0.65\textwidth]{./Figs9/ori_eigfun_spec_s_0} \\
\small
$e(k) = \mathcal{O}\left(k^{-1}\right)$ implies $L^2$ summability;
$e(k) = \mathcal{O}\left(k^{-1.65}\right)$ means \blue{$\phi_+ \in H^{0.28}(\TT^2)$}

     
\Emp
}

%%%%%%%%%%%%%%%%%%%%%%%%%%%%%%%%%%%%%%%%%%%%%%%%%%%%%%%%%%%%%%%%%%%%%%%%%%
%%%%%%%%%%%%%%%%%%%%%%%%%%%%%%%%%%%%%%%%%%%%%%%%%%%%%%%%%%%%%%%%%%%%%%%%%%
%%%%%%%%%%%%%%%%%%%%%%%%%%%%%%%%%%%%%%%%%%%%%%%%%%%%%%%%%%%%%%%%%%%%%%%%%%
\subsection{Exponential Growth}
\frame
{
\hspace*{-0.75cm}
\Bmp{12cm}
\raggedright

Exponential growth of the perturbation
\bigskip

\centering
\mbox{
  \includegraphics[width=0.475\textwidth]{./Figs9/eig_linear_growth} \quad
  \includegraphics[width=0.475\textwidth]{./Figs9/eig_linear_growthrate}}
\\
\medskip
The actual growth rate is within 0.01\% equal $\Re(\lambda_+^{3000})$

\Emp
}


%%%%%%%%%%%%%%%%%%%%%%%%%%%%%%%%%%%%%%%%%%%%%%%%%%%%%%%%%%%%%%%%%%%%%%%%%%
%%%%%%%%%%%%%%%%%%%%%%%%%%%%%%%%%%%%%%%%%%%%%%%%%%%%%%%%%%%%%%%%%%%%%%%%%%
%%%%%%%%%%%%%%%%%%%%%%%%%%%%%%%%%%%%%%%%%%%%%%%%%%%%%%%%%%%%%%%%%%%%%%%%%%
\subsection{Nonlinear Instability}
\frame
{
\hspace*{-0.75cm}
\Bmp{12cm}
\raggedright
\centering

\mbox{
\includegraphics[width=0.45\textwidth]{./Figs9/growth_nonlinear}\qquad
\includegraphics[width=0.45\textwidth]{./Figs9/growth_rate_nonlinear}
}
\small
\medskip

The time-dependence of (left) the norm
{$\| \omega^{\delta}(t) - \omega_s\|_{H^{-1}}$} and (right) the growth
rate $(d/dt) \ln (\| \omega^{\delta}(t) - \omega_s\|_{H^{-1}})$ in the
solution of the nonlinear problem with the initial condition given in
terms of the eigenfunction $\phi_+^{200}$ as
$\omega_0 = \omega_s + \delta \cdot
\Re\left(\phi_+^{200}\right)/||\Re\left(\phi_+^{200}\right)||_{H^{-1}}$
with different indicated magnitudes $\delta$.

\Emp
}


%%%%%%%%%%%%%%%%%%%%%%%%%%%%%%%%%%%%%%%%%%%%%%%%%%%%%%%%%%%%%%%%%%%%%%%%%%
%%%%%%%%%%%%%%%%%%%%%%%%%%%%%%%%%%%%%%%%%%%%%%%%%%%%%%%%%%%%%%%%%%%%%%%%%%
%%%%%%%%%%%%%%%%%%%%%%%%%%%%%%%%%%%%%%%%%%%%%%%%%%%%%%%%%%%%%%%%%%%%%%%%%%
\section{Nonmodal Growth of Perturbations}
\subsection{Optimization Formulation}
\frame
{
\hspace*{-0.75cm}
\Bmp{12cm}
\raggedright

\begin{itemize}
\item<1-> The growth rate of the eigenfunction with $\Re(\lambda_+0)
  \approx 0.1424$ does not saturate the bound on the growth of the
  semigroup
\begin{equation*}
\blue{ \left \| e^{\L t} w_0 \right \|_{H^{m}} \le \left \| w_0 \right \|_{H^{m}} e^{\mumax t}, \qquad \mumax = 1, \quad m = \pm 1}
\end{equation*}

\bigskip
\item<2-> How to find initial data $w_0$ realizing this growth?

 \bigskip
\item<3->
  Given $T > 0$, define the objective functional $J_m \; : \;
  {H_0^m(\TT^2)} \rightarrow \RR$
\begin{equation*}
\blue{  J_m (w_0) = \left|\left|e^{T\L} w_0 \right|\right|^2_{H^m},
\qquad m = \pm 1}
\end{equation*}
  

\end{itemize}
\Emp
}

%%%%%%%%%%%%%%%%%%%%%%%%%%%%%%%%%%%%%%%%%%%%%%%%%%%%%%%%%%%%%%%%%%%%%%%%%%
%%%%%%%%%%%%%%%%%%%%%%%%%%%%%%%%%%%%%%%%%%%%%%%%%%%%%%%%%%%%%%%%%%%%%%%%%%
%%%%%%%%%%%%%%%%%%%%%%%%%%%%%%%%%%%%%%%%%%%%%%%%%%%%%%%%%%%%%%%%%%%%%%%%%%
\frame
{
\hspace*{-0.75cm}
\Bmp{12cm}
\raggedright

\begin{itemize}
\item<1-> []
\begin{problem}[Optimization]
\vspace*{-0.5cm}
\blue{\small
\begin{align*}
& \max_{w_0 \in \mc M} J_m (w_0), \quad \text{where}  \ m = \pm 1  \ \text{and} \\
& \qquad \mc M := \{w_0 \in H^m_0(\TT^2):  \|w_0\|_{H^m} =   1\}. \\ \\
& \text{subject to:} \quad  \left\{
\begin{aligned}
&\partial_tw - \L w = 0,  && \text{in} \ \TT^2\times(0,T] \\ 
&w = w_0                 &&  \text{in} \ \TT^2 \ \text{at} \ t=0 \\
&\textrm{Periodic Boundary Conditions}                  &&  
\end{aligned}
%\label{eq:NS}
\right.
\end{align*}
}
\end{problem}

\item<2->
  Riemannian optimization problem \\ \qquad $\Longrightarrow$ \quad solution with a discretized gradient flow
\end{itemize}
\Emp
}

%%%%%%%%%%%%%%%%%%%%%%%%%%%%%%%%%%%%%%%%%%%%%%%%%%%%%%%%%%%%%%%%%%%%%%%%%%
%%%%%%%%%%%%%%%%%%%%%%%%%%%%%%%%%%%%%%%%%%%%%%%%%%%%%%%%%%%%%%%%%%%%%%%%%%
%%%%%%%%%%%%%%%%%%%%%%%%%%%%%%%%%%%%%%%%%%%%%%%%%%%%%%%%%%%%%%%%%%%%%%%%%%
\frame
{
\hspace*{-0.75cm}
\Bmp{12cm}
\raggedright

\begin{itemize}
\item<1-> Riemannian conjugate-gradient method: \quad $\widetilde{w}_{0} = \lim_{n \rightarrow \infty} w^{(n)} _0$
 \begin{equation*} 
\red{w^{(n+1)}_0  = \R\Big[w^{(n)} _0+ \tau_n d^{(n)}\Big]}, \qquad n=0,1,\dots,
\end{equation*}
where
\begin{itemize}
\item<2-> retraction $\R \; : \; H^{m} \rightarrow \M$, \qquad  $\R (w) := \frac{w}{||w||_{H^{m}}}$, \quad $m=\pm 1$

\smallskip
\item<3->{}
descent direction  \quad $d^{(n)} = \P_n  \blue{\nabla J\left(w^{(n)}_0\right)}+ \beta_n \Gamma_{\tau_{n-1} d^{(n-1)}}\left({d^{(n-1)}}\right)${}

\smallskip
\item<4->{}
gradient \quad \blue{$\nabla J(w_0) = 2 (1+|D|^2) e^{T\L^*} (1+|D|^2)^{-1} e^{T\L} w_0$} \\
obtained by solving the adjoint system (a terminal-value problem)
\blue{
\begin{align*}
- \partial_t w & = \L^* w, \qquad \L^*w = \u_s\cdot \nabla w + |D|^{-2}(\nabla w_s \cdot \nabla^\perp w), \\
\qquad w|_{t = T} & = (1+|D|^2)^{-1} e^{T\L} w_0
\end{align*}}
\vspace*{-0.5cm}

\item<5-> orthogonal projection onto tangent subspace \quad $\P_n \; : \; H^{m} \rightarrow \T_{w^{(n)} _0}\M${}

\smallskip
\item<6-> vector transport \qquad 
$\Gamma_{\tau_{n-1} d^{(n-1)}} \; : \; \T_{w^{(n-1)} _0}\M  \rightarrow \T_{w^{(n)} _0}\M$
  

\end{itemize}

  

\end{itemize}
\Emp
}

%%%%%%%%%%%%%%%%%%%%%%%%%%%%%%%%%%%%%%%%%%%%%%%%%%%%%%%%%%%%%%%%%%%%%%%%%%
%%%%%%%%%%%%%%%%%%%%%%%%%%%%%%%%%%%%%%%%%%%%%%%%%%%%%%%%%%%%%%%%%%%%%%%%%%
%%%%%%%%%%%%%%%%%%%%%%%%%%%%%%%%%%%%%%%%%%%%%%%%%%%%%%%%%%%%%%%%%%%%%%%%%%
\frame
{
\hspace*{-0.75cm}
\Bmp{12cm}
\raggedright
\frametitle{Riemannian Conjugate-Gradient Method \\ (Absil et al., 2008)}
\centering{}

\includegraphics[width=1.0\textwidth]{./Figs9/rcg}

\Emp
}



%%%%%%%%%%%%%%%%%%%%%%%%%%%%%%%%%%%%%%%%%%%%%%%%%%%%%%%%%%%%%%%%%%%%%%%%%%
%%%%%%%%%%%%%%%%%%%%%%%%%%%%%%%%%%%%%%%%%%%%%%%%%%%%%%%%%%%%%%%%%%%%%%%%%%
%%%%%%%%%%%%%%%%%%%%%%%%%%%%%%%%%%%%%%%%%%%%%%%%%%%%%%%%%%%%%%%%%%%%%%%%%%
\subsection{Flows Saturating the Spectral Bound}
\frame
{
\hspace*{-0.75cm}
\Bmp{12cm}
\raggedright
\centering

\begin{figure}
\mbox{\subfigure[$N=128$]{\includegraphics[width=0.325\textwidth]{./Figs9/s_1_128}}\qquad
\subfigure[$N=256$]{\includegraphics[width=0.325\textwidth]{./Figs9/s_1_256}}}
\mbox{\subfigure[$N=512$]{\includegraphics[width=0.325\textwidth]{./Figs9/s_1_512}}\qquad
\subfigure[$N=1024$]{\includegraphics[width=0.325\textwidth]{./Figs9/s_1_1024}}}
\end{figure}

\small\vspace*{-0.3cm}
\blue{[$m=1$]}  The optimal initial conditions $\widetilde{w}_0^N$ obtained using
different resolutions $N$.

\Emp
}


%%%%%%%%%%%%%%%%%%%%%%%%%%%%%%%%%%%%%%%%%%%%%%%%%%%%%%%%%%%%%%%%%%%%%%%%%%
%%%%%%%%%%%%%%%%%%%%%%%%%%%%%%%%%%%%%%%%%%%%%%%%%%%%%%%%%%%%%%%%%%%%%%%%%%
%%%%%%%%%%%%%%%%%%%%%%%%%%%%%%%%%%%%%%%%%%%%%%%%%%%%%%%%%%%%%%%%%%%%%%%%%%
\frame
{
\hspace*{-0.75cm}
\Bmp{12cm}
\raggedright

\mbox{
\includegraphics[width=0.475\textwidth]{./Figs9/s_1_growth}
\includegraphics[width=0.475\textwidth]{./Figs9/approx_s_1_crlt}}
\centering
\small
\blue{[$m=1$]} (left) The growth rate
  $(d/dt) \ln (\|w^N(t)\|_{H^{1}})$ versus $t$ for
  the optimal initial conditions $\widetilde{w}_0^N$ obtained using increasing
  spatial resolutions {$N^2$}. \\
  (right) The time-dependence of the
  autocorrelation function
  \begin{equation*}
\C^1_{\tau}(t) := \frac{\left\langle w(\tau), w(t) \right\rangle_{H^1}}{\| w(\tau) \|_{H^1} \| w(t) \|_{H^1}}, 
\qquad t, \tau \ge 0
\end{equation*}
  corresponding to
  $N = 1024$ and $\tau = 0, 0.25, 0.5, 0.75, 1$. \\
  Similar results obtained for $m = -1$.
\Emp
}



%%%%%%%%%%%%%%%%%%%%%%%%%%%%%%%%%%%%%%%%%%%%%%%%%%%%%%%%%%%%%%%%%%%%%%%%%%
%%%%%%%%%%%%%%%%%%%%%%%%%%%%%%%%%%%%%%%%%%%%%%%%%%%%%%%%%%%%%%%%%%%%%%%%%%
%%%%%%%%%%%%%%%%%%%%%%%%%%%%%%%%%%%%%%%%%%%%%%%%%%%%%%%%%%%%%%%%%%%%%%%%%%
\frame
{
\hspace*{-0.75cm}
\Bmp{12cm}
\raggedright
\centering

\movie[width=0.9\textwidth,loop,showcontrols]{\includegraphics[width=0.9\textwidth]{./Figs9/s_1_1024}}{../Euler2D_stability/Movies/Movie1.mov}

%\movie[width=0.8\textwidth,loop,showcontrols]{\includegraphics[width=0.8\textwidth]{./Figs9/s_1_1024}}{../Euler2D_stability/Movies/movie1.mp4 }

\Emp
}


%%%%%%%%%%%%%%%%%%%%%%%%%%%%%%%%%%%%%%%%%%%%%%%%%%%%%%%%%%%%%%%%%%%%%%%%%%
%%%%%%%%%%%%%%%%%%%%%%%%%%%%%%%%%%%%%%%%%%%%%%%%%%%%%%%%%%%%%%%%%%%%%%%%%%
%%%%%%%%%%%%%%%%%%%%%%%%%%%%%%%%%%%%%%%%%%%%%%%%%%%%%%%%%%%%%%%%%%%%%%%%%%
\frame
{
\frametitle{Conclusions}
\vspace{-0.1cm}
%\vspace*{0.0cm}
\hspace*{-1.0cm}
\Bmp{12cm}
\raggedright
\small
\begin{itemize}
\item<1-> The Taylor-Green vortex is linearly unstable
  \begin{itemize}
\item<2-> the unstable eigenvalue embedded in the essential spectrum, \\ \hfill \blue{$\lambda_{+} \in \sigma_{\text{ess}}(\L)$}
\medskip
\item<3-> the corresponding eigenfunction given by a distribution with the Sobolev
  regularity consistent with Lin (2004)
  \medskip 
\item<4-> the equilibrium also exhibits a nonmodal growth mechanism involving a
  continuous family of uncorrelated functions which saturates the
  spectral bound
 \begin{itemize}
 \item<5-> intrinsically related to the $\infty$-dimensional nature of the operator $\L$ 
 \smallskip
 \item<6-> the mechanism is consistent with the WKB analysis based on
   the bicharacteristic equation
\end{itemize}
\end{itemize}

  \bigskip
\item<7-> \blue{\sc Open Question:} stability of equilibria in 3D
\end{itemize}
\Emp
}



\end{document}

%%%%%%%%%%%%%%%%%%%%%%%%%%%%%%%%%%%%%%%%%%%%%%%%%%%%%%%%%%%%%%%%%%%%%%%%%%
%%%%%%%%%%%%%%%%%%%%%%%%%%%%%%%%%%%%%%%%%%%%%%%%%%%%%%%%%%%%%%%%%%%%%%%%%%
%%%%%%%%%%%%%%%%%%%%%%%%%%%%%%%%%%%%%%%%%%%%%%%%%%%%%%%%%%%%%%%%%%%%%%%%%%

\frame{
  \hspace*{-1.0cm}
 \frametitle{References}
\Bmp{12cm}
\raggedright
%\centering
%\vspace{-3.5cm}
\begin{itemize}
\item<1-> A.~Elcrat and B.~Protas, ``A Framework for Linear Stability
  Analysis of Finite-Area Vortices'', {\em Proceedings of the Royal
    Society A}, {\bf{469}}, 20120709, 2013.

 \bigskip
\item<1-> B.~Protas and A.~Elcrat, ``Linear Stability of Hill's Vortex
  to Axisymmetric Perturbations'', {\em Journal of Fluid Mechanics}
  {\bf{799}}, 579--602, 2016.

 \bigskip
\item<1-> B.~Protas, ``Linear Stability of Inviscid Vortex Rings to
  Axisymmetric Perturbations'', {\em Journal of Fluid Mechanics},
  {\bf{874}}, 1115--1146, 2019.

\end{itemize}

\Emp
}




